% GENERATED, edit the French unit or the translation table
\mxunit{premiere-c05}{Conditional probability and independence}\label{C05}
\begin{keypoint}[Chapter path]
\textbf{Objectives.} Change the reference sample space; build a tree; use a partition; distinguish disjoint events from independent events; handle up to four Bernoulli trials.

\textbf{Prerequisites.} Fractions, proportions, probabilities on a finite sample space, intersection and union.

\textbf{Yearly reference point.} 17--26 November: 8 hours\end{keypoint}
\section*{Diagnostic: \mxcount{diag} questions, 10 minutes}\label{C05-D}
Answer without consulting the lesson. Mark the questions for which help is needed. This diagnostic is not graded.
\begin{enumerate}
\item Out of 40 students, 12 play a given sport. What proportion is that?
\item If $P(A)=\num{0,3}$, give $P(\overline A)$.
\item Write ``A and B'' in set notation.
\item Write ``A or B'' using the inclusive or.
\item Calculate $\num{0,4}\times\num{0,25}$.
\item A fair die is rolled. Give $P(\ens{2;4;6})$.
\mxcounted{diag}
\end{enumerate}
\paragraph{After the diagnostic.} Revisit the flagged topics with the teacher; redo a similar question, then explain the rule used.
\paragraph{Speaking and logic.} Explain aloud what ``given'', ``and'' and ``at least'' mean. On a tree, justify an addition and a multiplication.
\mxnewpage
\section*{Investigate before formalising}\label{C05-RECH}
Build a table of 100 people in which A and B are independent, with 40 in A and 30 in B. Change two cells, keeping these totals, so as to destroy the independence.
\paragraph{Organisation.} Individual investigation 5 min, comparison in pairs 5 min, class discussion 5 min. Keep a record of one attempt, even an unfinished one. A conjecture must then be justified.
\mxnewpage
\section{Conditioning: changing the denominator}\label{C05-U1}
\begin{keypoint}[Definition and results]
For $P(A)>0$, the probability of B given A is $P_A(B)=\dfrac{P(A\cap B)}{P(A)}$. In a table of equally likely counts, it is the quotient of the count of $A\cap B$ by the count of A. In general $P_A(B)\ne P_B(A)$.
\end{keypoint}
\paragraph{Status and justification.} Definition. Multiplying by $P(A)$ gives the rule $P(A\cap B)=P(A)P_A(B)$. The condition $P(A)>0$ is essential.
\begin{example}Out of 100 people, 40 belong to A, 30 to B and 12 to both. $P_A(B)=\dfrac{12}{40}=\num{0,3}$, whereas $P_B(A)=\dfrac{12}{30}=\num{0,4}$. $P(A\cap B)=\num{0,12}$.\end{example}\label{C05-EX1}
\begin{pitfall}``Among the A's'' fixes the denominator. A conditional probability is not an intersection.\end{pitfall}
\begin{application}Out of 60 objects, 20 are red; among the red ones, 5 are damaged. Calculate $P(R\cap D)$ and $P_R(D)$.\end{application}\label{C05-AP1}
\section{Trees, partitions and total probability}\label{C05-U2}
\begin{keypoint}[Definition and results]
A partition $A_1,\ldots,A_k$ is a family of disjoint events whose union is the sample space. If their probabilities are positive, $P(B)=\sum_iP(A_i)P_{A_i}(B)$. On a tree, the probabilities leaving the same node sum to 1; the probabilities along a path are multiplied; disjoint paths are combined by adding.
\end{keypoint}
\paragraph{Status and justification.} $B$ is the disjoint union of the sets $B\cap A_i$. By additivity and by the definition of conditional probability, $P(B)=\sum_iP(B\cap A_i)=\sum_iP(A_i)P_{A_i}(B)$.
\begin{example}Two machines supply \pc{60} and \pc{40} of a production; their defect rates are \pc{2} and \pc{5}. $P(D)=\num{0,6}\times\num{0,02}+\num{0,4}\times\num{0,05}=\num{0,032}$. Among the defective items, the first machine accounts for $\dfrac{\num{0,012}}{\num{0,032}}=\dfrac{3}{8}$.\end{example}\label{C05-EX2}
\mxinclude{figures/C05-U2}
\begin{pitfall}The simple average of the rates is correct only if the branches have equal weights.\end{pitfall}
\begin{application}$P(A)=\dfrac{1}{3}$, $P_A(B)=\dfrac{3}{4}$, $P_{\overline A}(B)=\dfrac{1}{4}$. Calculate $P(B)$.\end{application}\label{C05-AP2}
\section{Independence and disjointness}\label{C05-U3}
\begin{keypoint}[Definition and results]
A and B are independent if $P(A\cap B)=P(A)P(B)$. If $P(A)>0$, this is equivalent to $P_A(B)=P(B)$. They are disjoint if $A\cap B=\varnothing$. Two disjoint events of strictly positive probability are not independent.
\end{keypoint}
\paragraph{Status and justification.} The conditional criterion follows from a division by $P(A)>0$. For disjoint events of positive probability, the left-hand side is 0 and the right-hand product is strictly positive.
\begin{example}With a fair die, let A: the result is even, B: the result is at least 5. $P(A)=\dfrac{1}{2}$, $P(B)=\dfrac{1}{3}$, $P(A\cap B)=\dfrac{1}{6}$: independent. A and ``the result is odd'' are disjoint but not independent.\end{example}\label{C05-EX3}
\begin{pitfall}Independence in a model must be justified or stated explicitly; ``two different events'' is not enough.\end{pitfall}
\begin{application}Given $P(A)=\num{0,4}$, $P(B)=\num{0,5}$, $P(A\cap B)=\num{0,18}$. Are they independent?\end{application}\label{C05-AP3}
\section{Successive trials and Bernoulli trials}\label{C05-U4}
\begin{keypoint}[Definition and results]
A Bernoulli trial has two outcomes: success, with probability p, and failure, with probability $1-p$. In a repetition of n independent, identical trials, with $n\leq4$ here, the branch probabilities remain p and $1-p$. A path with k successes has probability $p^k(1-p)^{n-k}$; we count the paths that realise the event.
\end{keypoint}
\paragraph{Status and justification.} The probability of a path is a product, by independence. The paths corresponding to an event are disjoint: their probabilities add.
\begin{example}For three trials with $p=\dfrac{1}{2}$, exactly two successes correspond to SSF, SFS, FSS. Each path has probability $\dfrac{1}{8}$, so the total is $\dfrac{3}{8}$. At least one success has probability $1-(\dfrac{1}{2})^3=\dfrac{7}{8}$.\end{example}\label{C05-EX4}
\begin{pitfall}When drawing without replacement from an urn, the probabilities generally change after the first draw: the independent model may be wrong.\end{pitfall}
\begin{application}Four independent trials, each with probability $p=\dfrac{1}{3}$. Calculate the probability of four failures, then of at least one success.\end{application}\label{C05-AP4}
\section{Simulating and discussing the model}\label{C05-U5}
\begin{keypoint}[Definition and results]
A pseudo-random simulation produces frequencies that fluctuate around a probability of the model. For a uniform point in the square $\intff{0}{1}^2$, the proportion satisfying $y\leq x^2$ estimates the area under the parabola; the proportion satisfying $x^2+y^2\leq1$, multiplied by 4, estimates $\pi$. Repetitions must use new draws.
\end{keypoint}
\paragraph{Status and justification.} Simulation principle accepted and explored experimentally. The equality of probability and area is an interpretation of the uniform draw; integration is not a prerequisite. A simulation does not prove an exact value.
\begin{example}Among 10000 points, if 7840 lie in the quarter disc, the estimate of $\pi$ is $\dfrac{4\times 7840}{10000}=\num{3,136}$. Two trials of the same size generally give different values.\end{example}\label{C05-EX5}
\begin{pitfall}Fixing a seed makes a trial reproducible; it neither removes the randomness of the model nor guarantees any accuracy.\end{pitfall}
\begin{application}A simulation of 2000 points gives 675 points below $y=x^2$. What area does it estimate?\end{application}\label{C05-AP5}
\mxnewpage
\section*{Exercise path}\label{C05-PARC}
Start with exercises 1 and 2. Then choose one of exercises 3 or 4 and do it without help. Exercises 5 to 7 are a reserve for transfer, investigation or programming: the teacher designates one. Exercise 8 is revisited after a delay.
\begingroup
\mxexercisehead{C05-E1}{A table}
A survey covers 100 users of two transport services A and B. Of these, 40 use A, 35 use B and 15 use both. One of the 100 users is chosen at random; all are equally likely to be chosen. The events $A$ and $B$ mean, respectively, ``the user uses service A'' and ``the user uses service B''.
\begin{enumerate}[label=\arabic*.]
\item Build and complete a two-way table: use or non-use of A in the rows, use or non-use of B in the columns. Add the totals.
\item Calculate $P_A(B)$ and $P_B(A)$. For each calculation, state the reference population used as the denominator.
\item Interpret the two probabilities in two distinct sentences.
\end{enumerate}
\mxexercisehead{C05-E2}{Production}
A company makes objects in two workshops A and B, which provide $\dfrac14$ and $\dfrac34$ of the production respectively. Among the objects made by A, 4 percent are defective; among those made by B, 2 percent are. An object is chosen at random from the whole production. Let $D$ be the event ``the object is defective''.
\begin{enumerate}[label=\arabic*.]
\item Represent the situation with a weighted tree and state the meaning of each second-level probability.
\item Calculate $P(D)$ and interpret the result as an overall defect rate.
\item It is now known that the chosen object is defective. What is the probability that it was made in workshop A?
\end{enumerate}
\mxexercisehead{C05-E3}{Checking independence}
A fair six-sided die, with faces numbered 1 to 6, is rolled. Consider the events $A$: ``obtain an even number'' and $B$: ``obtain a number greater than or equal to 4''.
\begin{enumerate}[label=\arabic*.]
\item Write the sets of outcomes corresponding to $A$, $B$ and $A\cap B$.
\item Determine whether the events $A$ and $B$ are independent. Justify using a numerical criterion.
\item Explain why they are not disjoint either.
\end{enumerate}
\mxexercisehead{C05-E4}{Without replacement}
An urn contains three red balls and two blue balls, indistinguishable by touch. A first ball is drawn at random, then a second without replacing the first. At each draw, all the balls in the urn are equally likely.
\begin{enumerate}[label=\arabic*.]
\item Build the weighted tree of the two draws, distinguishing the colour obtained at the first and at the second draw.
\item Calculate the probability of drawing two red balls.
\item Calculate the probability of drawing one ball of each colour. Explain why two paths of the tree must be taken into account.
\end{enumerate}
\mxexercisehead{C05-E5}{Three trials}
An experiment has a success of probability $\dfrac13$ and a failure. It is repeated three times independently, under the same conditions.
\begin{enumerate}[label=\arabic*.]
\item Build a three-level tree, or list the eight possible sequences of outcomes.
\item Calculate the probability of exactly one success. State the paths used.
\item Calculate the probability of at least two successes, treating the cases of two and of three successes separately.
\end{enumerate}
\mxexercisehead{C05-E6}{Detector}
A device checks objects of which 2 percent are defective. It raises an alarm for 90 percent of the defective objects and, wrongly, for 5 percent of the objects without a defect. An inspected object is chosen at random. Let $D$ be the event ``the object is defective'' and $S$ the event ``the device raises an alarm''.
\begin{enumerate}[label=\arabic*.]
\item Determine the proportion of truly defective objects among those flagged. Choose a representation and detail the reasoning.
\item Give an exact value, then a percentage rounded to the nearest tenth.
\item Someone proposes 90 percent directly as the answer. Explain precisely which confusion of probabilities they are making.
\end{enumerate}
\mxexercisehead{C05-E7}{Sets and cardinality}
We work in the set $E=\ens{1;2;3;4}$. Consider its subsets $A=\ens{1;2}$ and $B=\ens{2;4}$.
\begin{enumerate}[label=\arabic*.]
\item Write $A\cap B$, $A\cup B$ and the complement of $A$ in $E$. Give the cardinality of each.
\item Write all the pairs of $A\times B$, then give its cardinality.
\item Explain why the pairs $\pt{1;2}$ and $\pt{2;1}$ cannot be confused. Does the second belong to $A\times B$?
\end{enumerate}
\mxexercisehead{C05-E8}{Delayed review}
In a random experiment, two events $A$ and $B$ satisfy\[P(A)=\dfrac12,\qquad P_A(B)=\dfrac14,\qquad P_{\overline A}(B)=\dfrac34.\]
\begin{enumerate}[label=\arabic*.]
\item Calculate $P(B)$ using a tree or the law of total probability.
\item Calculate $P_B(A)$, making the common event used in the numerator explicit.
\item Are the probabilities $P_A(B)$ and $P_B(A)$ equal here? Is this equality a general property?
\end{enumerate}
\paragraph{Homework.} Do only the task designated at the end of the session, in the dated progression: 10--20 min at most. The exercises depend on the topics already studied, so some questions are deferred. Revisit exercise 8 one week after the chapter (10 min). If the allotted time is exceeded, note the difficulty and bring your attempts; do not pile up the whole reserve.
\endgroup
\mxneed{12}\section*{To do in the CQFD textbook}
These exercises are done in your textbook. The progression indicates the session; they replace the homework planned for that day (10 to 20 minutes in total).
\paragraph{CQFD, exercise 67 p.~336, questions 1, 2 and 3.}\label{C05-M1} Build a tree keeping the textbook's events D and N. Careful: N denotes a negative test.
\paragraph{CQFD, exercise 70 p.~336, questions a and b.}\label{C05-M2} Count the four equally likely outcomes of the two throws to refute the proposed equalities.
\begingroup\small
\mxnewpage\section*{MCQ 1: understand and choose}\label{C05-Q1}
\mxCount{qcm1} single-answer questions for practice. Choose one answer per question; then explain one distractor. This chapter format does not reproduce the official EAM format.
\begin{enumerate}
\item $P_A(B)$ has as denominator
\begin{enumerate}[label=\Alph*.,itemsep=2pt]
\item $P(B)$
\item $P(A)$
\item $P(A\cap B)$
\item $1-P(A)$
\end{enumerate}
\item $P(A)=\num{0,4}$, $P_A(B)=\num{0,2}$: intersection
\begin{enumerate}[label=\Alph*.,itemsep=2pt]
\item $\num{0,08}$
\item $\num{0,6}$
\item $\num{0,5}$
\item $\num{0,2}$
\end{enumerate}
\item A partition consists of events that are
\begin{enumerate}[label=\Alph*.,itemsep=2pt]
\item equal
\item all independent
\item equally probable
\item disjoint and covering the sample space
\end{enumerate}
\item Independence is equivalent to
\begin{enumerate}[label=\Alph*.,itemsep=2pt]
\item $P(A)=P(B)$
\item $P(A\cap B)=0$
\item $P(A\cap B)=P(A)P(B)$
\item $A=B$
\end{enumerate}
\item Two disjoint events of positive probability are
\begin{enumerate}[label=\Alph*.,itemsep=2pt]
\item independent
\item not independent
\item equal
\item necessarily complementary
\end{enumerate}
\mxcounted{qcm1}
\end{enumerate}
\endgroup
\begingroup\small
\mxneed{18}\section*{MCQ 2: understand and choose}\label{C05-Q2}
\mxCount{qcm2} single-answer questions for practice. Choose one answer per question; then explain one distractor. This chapter format does not reproduce the official EAM format.
\begin{enumerate}
\item Three independent tosses of a fair coin, at least one head:
\begin{enumerate}[label=\Alph*.,itemsep=2pt]
\item $\dfrac{1}{8}$
\item $\dfrac{3}{8}$
\item $\dfrac{7}{8}$
\item $\dfrac{1}{2}$
\end{enumerate}
\item Two independent trials with $p=\dfrac{1}{3}$, two successes:
\begin{enumerate}[label=\Alph*.,itemsep=2pt]
\item $\dfrac{1}{9}$
\item $\dfrac{2}{3}$
\item $\dfrac{1}{3}$
\item $\dfrac{4}{9}$
\end{enumerate}
\item The branches leaving the same node sum to
\begin{enumerate}[label=\Alph*.,itemsep=2pt]
\item $0$
\item $2$
\item the probability of the path
\item $1$
\end{enumerate}
\item A simulation gives a frequency equal to \num{0,51}. This
\begin{enumerate}[label=\Alph*.,itemsep=2pt]
\item proves $p=\num{0,51}$
\item is a fluctuating observation
\item proves $p=\dfrac{1}{2}$
\item refutes $p=\dfrac{1}{2}$
\end{enumerate}
\item Drawing without replacement, after a red ball from 3 red and 2 blue, the probability of red is:
\begin{enumerate}[label=\Alph*.,itemsep=2pt]
\item $\dfrac{3}{5}$
\item $\dfrac{3}{4}$
\item $\dfrac{1}{2}$
\item $\dfrac{2}{5}$
\end{enumerate}
\mxcounted{qcm2}
\end{enumerate}
\subsection*{True or false? Justify.}\label{C05-VF}
\begin{enumerate}
\item $P_A(B)=P_B(A)$ for all events of positive probability.
\item If A and B are independent, so are A and $\overline B$.
\item Two successive draws are always independent.
\end{enumerate}
\begin{keypoint}[Self-assessment]I identify the right denominator; I check the sums at the nodes; I justify independence; I count the paths without confusing them.

For each objective, indicate: achieved without help; achieved with help; to be reworked. Write one example and one mistake to avoid.\end{keypoint}
\endgroup
\ifmxteacher\mxinclude{corriges/C05}\fi
